\documentclass{ltjsbook}
\usepackage{docmute} 
\usepackage{../../myStyle} 

\begin{document}
% \chapter{心理尺度を作る}
\chapter{Creating a Psychological Scale}
\label{x02_testTheory}
% \section{はじめに}
\section{Introduction}
% 心理学の研究法は調査，実験，観察の3つが代表的なものです。とくに調査法は，質問紙を用いて多くの人に回答を依頼し，それを統計的に分析することで研究仮説を確認しようとするものです。調査対象者をどのように集めるか，調査票をどのようにデザインするか，得られたデータをどのように分析するか，といったことでそれぞれ90分以上話せるぐらいに色々考えるべきことはありますが，ここではとくにその本質について考えてみましょう。すなわち，なぜ紙で用紙に丸をつけたものを集めただけで，心の何かがわかったと思えるのか，ということです。
The three representative methods of psychological research are surveys, experiments, and observations. In particular, the survey method involves soliciting responses from many people using a questionnaire, and then statistically analyzing those responses to confirm research hypotheses. There are many things to consider such as how to gather the survey respondents, how to design the questionnaire, and how to analyze the obtained data, which can each take more than 90 minutes to discuss, but let's think specifically about the essence of it. That is, why do we think we understand something about the mind just by collecting marked papers on a form?

% アンケート調査は一見すると，なんのひねりもない研究法に思えます。すなわち，聞きたいことを当の本人に聞いてみる，これだけです。しかも答えやすいように(集計しやすいように？ )紙に問題が書いてあって，該当するところに丸をつけるだけでよかったりするのです。この手の研究は少なくとも100人，できれば数百人，大きい規模だと千以上の桁数の協力者に回答を求めます。それをPCに入力して集計するわけです\footnote{もっとも最近では，webで調査の回答を得ることで，入力や誤入力のチェックなどの手間が大幅に軽減されています。}。しかし「そう思う」を5点，「ややそう思う」を4点，以下3，2，1点・・・と入力するのはなぜでしょう。これらの反応は\keytermL{順序尺度水準}{じゅんじょしゃくどすいじゅん}でしかない，と言われますが，それを間隔尺度水準であると「みなし」て，平均値を求めたりします。なんでそんなことが許されるのでしょう。
At first glance, a survey seems like a straightforward research method. That is, the method simply involves asking the individuals in question what one wishes to know. Moreover, the questions are written on paper for easy answering (for easy tallying?). All that's required is to circle the response that applies. This type of research generally requests responses from at least a hundred people, preferably hundreds, and for large scale surveys, to the tune of a thousand or more. These responses are then entered into a computer for tallying\footnote{More recently, gathering survey responses via the web has greatly reduced the time and effort spent on tasks like data entry and error checking.}. But why is "strongly agree" assigned 5 points, "somewhat agree" 4 points, and so on down to 1 point...? These responses are said to be at the "ordinal scale level," but they are 'deemed' to be at the interval scale level, and averages are calculated using them. But why is such a practice allowed?

% 本稿ではその謎について迫っていきたいと思います。
In this article, I would like to delve into that mystery.
% \subsection{測ろうとしているもの；態度}
\subsection{What we are trying to measure; Attitude}
% 調査票で質問する内容は，(お客様アンケートとかではなく)心理学の研究であれば\keyterm{尺度}{しゃくど}{scale}を使います。尺度とは，測定したい対象に数字を割り振るルールのことであり，心理尺度は項目とその採点方法が規格化されたものを指します。心理尺度の測ろうとしているものの多くは\keyterm{態度}{たいど}{Attitude}と呼ばれるものです。これは社会心理学の用語で，簡単に定義すると態度とは「行動の準備状態」のことです\footnote{この定義は\textcite{Allport1935}によるものです。}。
The content of the questions on the survey form, if it's for psychological research (not customer surveys), uses what we call a "scale". A scale is a rule for assigning numbers to the target you want to measure, and a psychological scale refers to a system where the scoring method for each item is standardized. What most psychological scales try to measure is something called "Attitude". This is a term from social psychology, and when simply defined, attitude refers to the "state of preparedness for action". This definition is from Allport's 1935 study.

% 社会心理学も心理学ですから，研究は基本的に観察可能で客観的なものを対象にします。社会的な行動をテーマにするというのがもちろんですが，社会的な行動というのは状況によって変わるものですし，研究したい状況を待っていても自動的に出てくるものではありません。もちろん実験室などでその状況を作り出すこともやりますが，選挙のような公的なものであれば実験的に作り出すこともできません。しかしたとえば選挙結果などは，社会心理学における政治的行動に大きな影響を及ぼす(あるいは政治的行動の結果になる)ものですから，研究としては非常に重要なテーマになります。選挙になるまでは仕事ができない，というのでは社会心理学者も困りますから，「あなたはどこに投票するつもりですか」と聞いて普段の政治的な態度を調べるという研究手法ができました。社会調査の始まりです。
As social psychology is a part of psychology, its research primarily targets things that can be objectively observed. Of course, its theme is social behavior, which changes from situation to situation and doesn't automatically appear just by waiting for the desired situation to study. While we do create such situations in laboratories, we can't artificially construct public matters like elections. However, things like election results greatly influence political behavior in social psychology (or become the outcome of political behavior), so they are extremely important research topics. Social psychologists would be at a loss if they couldn't work until an election, so a research method was created where they would ask, "Where do you plan to vote?" to investigate usual political attitudes. This marked the beginning of social surveys.

% そしてこれを心理学一般に応用しているのが，質問紙調査です。心理学の研究テーマは行動ですが，行動を作り出せない場合や普段の状態を査定したいときに，「調査票に回答を求める」という方法を取るのです。もちろん調査に対する回答は，嘘や見栄，間違いや勘違いなども入り込む可能性がありますので，設計は丁寧に，分析は慎重に行う必要があります。調査票で過去のことを聞いたりすると記憶が歪んでいる可能性がありますし，未来のことを聞いたら嘘八百を答えられるかもしれません。今の気持ちを聞いても，自分の現在の状況がはっきりわかっているかどうか，怪しいものです。それでもある程度の真実味はあるだろうと考えて，こうした阻害要因を取り除いて本質を掴むために，いろいろな工夫がなされています。
The questionnaire survey is a broad application in psychology. Although the research theme of psychology is behavior, when it is not possible to create behavior or when it is desired to assess the usual condition, a method is adopted where responses are asked for in the questionnaire. Of course, responses to surveys may include lies, pretensions, mistakes, misunderstandings, etc., so the design must be carefully done, and the analysis must be done cautiously. If you ask about the past in the questionnaire, your memory might be distorted, and if you ask about the future, you might get absurd answers. Even if you ask about your current feelings, it's dubious whether you clearly understand your current situation. Still, we assume that there is some degree of truth. Various attempts have been made to remove these inhibitory factors and grasp the essence.

% 尺度を使って感情や気分，考え方(認知傾向)や今後どう行動するつもりか(行動意図)を調べることがなされていますが，本質的にはこれらすべてが，広い意味での態度です。態度の認知的側面，感情的側面，行動的側面などと言われたりします\parencite{Fujihara2001}。
Measurements are used to investigate emotions, moods, ways of thinking (cognitive tendencies), and future behavior intentions, but essentially, all of these are attitudes in a broad sense. They are often referred to as the cognitive aspect of attitude, the emotional aspect, and the behavioral aspect \parencite{Fujihara2001}.
% この「態度」の説明や定義は，実はそれほど明確ではありません。しかし，一般に特定の対象に対する，正負・量的な評価が可能なものと考えられています。たとえば「自民党に対する態度」というのは，自民党という対象に対して，「好き」「嫌い」というポジティブ・ネガティブの評価ができ，さらに「とても好き」「やや嫌い」のように量的に表現できるものでもある，と仮定されます。多くの人に多くの項目で調査するのは，項目同士の誤差が相殺しあってこれが正規分布すると考えられるからであり，社会的な態度は極端な値が少なく平均的な値を持つひとが多いものとして，相対的に評価されます。
The explanation or definition of this "attitude" is not really clear. However, it is generally considered to be something that can be evaluated positively or negatively, or quantitatively, towards a specific object. For example, the term "attitude towards the Liberal Democratic Party" assumes that it is possible to make positive and negative evaluations such as "like" and "dislike" towards the object called the Liberal Democratic Party, and further expressed quantitatively as "really like" and "slightly dislike". The reason for surveying many people on many items is because it is assumed that the errors between items will cancel each other out and this will be normally distributed, and social attitudes are relatively evaluated with fewer extreme values and many people holding average values.

% 基本的に測定しようとしているものは，こうした特徴を持ったなにかである，ということをまずは踏まえておきたいと思います\footnote{性格のように特定の対象を持たないものであっても，正規分布に従うと考えるのは自然ですから，この後説明する統計技法が適用できるものもあります。感情や気分といったものは，持続時間が短いので生理指標などで測定するべきであり，調査票によるアプローチは不向きですが，逆に言えばある程度一定の安定した心理状態であれば測定することができると考えられているのかもしれません。}。
The first thing to keep in mind is that what we are trying to measure is something with these features\footnote{Even if it doesn't have a specific target like a character, it's natural to assume that it follows a normal distribution, so some of the statistical techniques explained later can be applied. Emotions or moods, which have a short duration, should be measured with physiological indicators, and the questionnaire approach is not suitable. On the other hand, it may be thought that if it is a certain stable psychological state, it can be measured.}.

% \subsection{3つの方法}
\subsection{Three Methods}
% 心理尺度の作り方には大きく分けて3つのスタイルがあります。
There are roughly three styles of creating psychological scales.

% 1つ目はサーストン法による尺度で，\keyterm{等現間隔法}{とうげんかんかくほう}{method of equal-appearing intervals}と呼ばれるものです。2つ目はリッカート法(Likert法，ライカートと読む人もいる)と呼ばれるもので，5件法，7件法など数段階のカテゴリラベルのもっとも近いところに丸をつけるという方法です。3つ目はSD(Semantic Differential法，意味微分法と訳されることも)法とよばれるものです。SD法は態度測定というより，イメージの測定を目的としたもので，対象を提示しつつペアになった形容詞を列挙して提示します。たとえば「専修大学」という対象に対して，「激しい-落ち着いた」「慎重な-軽快な」といった形容詞対ではどちらの表現が近いかを評定してもらいます。形容詞の対が作る軸上でいうと平均的にどのあたりに対象がプロットされるのか，を見ることで対象ごとのイメージの違いを表現するのが基本的なアイデアです。SD法は複数の対象に対するイメージの相対的比較ですから，スコアの点数化にはそれほど重きを置いていないのでここでは取り上げません。
The first is a scale by the Thurstone method, which is called the "method of equal-appearing intervals". The second one is called the Likert method (some people pronounce it as "Lykaat"), which is a method where you put a circle at the closest place to the categorial label of several stages such as 5-point scale or 7-point scale. The third is called the SD (Semantic Differential) method. Rather than measuring attitudes, the SD method aims to measure images, presenting the target and listing paired adjectives. For example, for the target "Senshu University", the respondents rate which expression is closer in pairs of adjectives such as "intense-calm" or "careful-breezy". The basic idea is to express the difference in images for each target by seeing where the target is plotted on average on the axis created by the pair of adjectives. Since the SD method is a relative comparison of images for multiple targets, it does not give much importance to the quantification of scores, so we will not deal with it here.

% サーストン法とリッカート法は，これを使って態度のスコアをつけることができます。小杉の自民党に対する態度は4.8点だ，といったように，です。こうした数値化がどのような理屈でなされるのかを，今からみていこうと思います。
The Thurstone and Likert methods can be used to score attitudes. For example, Kokusai's attitude towards the Liberal Democratic Party is 4.8 points. Let's look at how such quantification is done.


% \section{サーストンの等現間隔法}
\section{Thurstone's Method of Equal Appearing Intervals}
\label{ThurstoneScale}
% サーストンの\ruby{等現間隔法}{とうげんかんかくほう}は，社会的な態度について絶対評価を与える方法です。この方法で作成された尺度は，各項目(態度表明文と呼ばれます)に尺度値がついており，回答者は提示された項目に賛成であればその尺度値がその人の態度得点になります。この尺度値は事前に複数の評定者によって決めておく必要があります。すなわち，尺度を作る前の入念な準備が必要です。また，サーストンの尺度は1次元性を有していることが前提となります。
Thurstone's Method of Equal Appearing Intervals is a method of giving absolute evaluations about social attitudes. The scales created by this method have scale values for each item (referred to as attitudinal statements) and if the respondent agrees with the presented item, its scale value becomes the person's attitude score. These scale values need to be predetermined by multiple evaluators in advance. In other words, careful preparation is a necessity before creating the scale. Also, Thurstone's scale requires the premise of having one-dimensionality.

% 具体的な作成方法は次のような手順で行います。
The specific creation method is performed in the following steps.
\begin{enumerate}
	\item 項目の収集
	\item 評定者集団による評定
	\item 尺度値の算出
	\item 項目の選定
\end{enumerate}


% 以下順に説明します。
I will explain the following in order.
% \paragraph{項目の収集}まずは測定したい社会的態度のテーマに沿って，項目を準備します。たとえば「自民党に対する態度」のように，誰でも思い描ける具体的な対象が良いでしょう。このようなテーマが決まれば，これに対する態度項目を色々考えます。「自民党のことを考えると夜も眠れない」とか「自民党に関係したニュースはなるべく見るようにしている」「近所の自民党員の事務所に行くことがある」というポジティブな態度もあるでしょうし，「自民党のニュースはなるべく聞きたくない」「自民党には投票しない」「自民党は不正まみれの悪い政党である」といったネガティブな態度もあるでしょう。こうした文言をなるべく多く，強い態度から弱い態度まで，ポジティブなものからネガティブなものまで網羅的に準備します。ニュートラルな項目も考えておく必要があります。
\paragraph{Item Collection} First, prepare the items along the theme of the social attitude you want to measure. For example, a concrete target that anyone can imagine, like "attitude towards the Liberal Democratic Party" would be good. Once such a theme is decided, various attitude items towards this are considered. There may be positive attitudes such as "I can't sleep at night thinking about the Liberal Democratic Party" or "I try to watch news related to the Liberal Democratic Party as much as possible", or "I visit the local office of the Liberal Democratic Party member". There may also be negative attitudes such as "I do not want to listen to the news about the Liberal Democratic Party as much as possible", "I will not vote for the Liberal Democratic Party", or "The Liberal Democratic Party is a corrupt and bad political party". Prepare as many of such expressions as possible, covering from strong to weak attitudes, and from positive to negative. It is also necessary to consider neutral items.

% \paragraph{評定者集団による評定}尺度値を決めるための事前準備です。まず評定者を無作為に集めます。少なくとも十数人は必要でしょう。評定者には事前に準備した項目が好意的－非好意的（または肯定的－否定的）の1次元にそって，7〜11段階ぐらいの多段階に分類してもらいます。「自民党のことを考えると夜も眠れない」というのは非常にポジティブなので11点，「自民党には投票しない」というのはかなりネガティブなので2点，といったようにです。
\paragraph{Rating by a Panel of Judges} This is a preparatory step to determining scale values. First, we gather a group of judges at random. At the very least, a dozen or so would be needed. The judges are tasked with categorizing pre-selected items along a single axis of favorability-unfavorability (or affirmation-negation), into multiple stages of about 7-11 steps. For example, if one was thinking about the Liberal Democratic Party and they couldn't sleep at night, this would be rated as very positive, thus given 11 points. On the other hand, if one said that they wouldn't vote for the Liberal Democratic Party, this would be considered quite negative, thus allocated 2 points.

% \paragraph{尺度値の算出}このように各態度表明文を複数人で評価してもらいますから，その項目の平均値，中央値，分散などの記述統計量を計算できます。この中央値(あるいは平均値)をその項目の尺度値とします。ただし，ここで分散が大きい項目は，評定者によって評定の仕方がバラバラだということを意味しますよね。値が人によって定まらないというのは，その項目が刺激としてあまり好ましくないと考えられるので，項目候補から削除します。誰がみても10点とか誰がみても3点，といった分散が少ない項目が望ましいでしょう。
\paragraph{Calculation of Scale Values} Since multiple people will evaluate each attitudinal statement, you can calculate descriptive statistics such as the mean, median, and variance of that item. The median (or mean value) is used as the scale value for that item. However, items with a large variance suggest that the method of rating varies by the rater. Items that do not have a consistent value across individuals are considered undesirable stimuli and are therefore removed from the list of candidate items. Items with little variance, like an item that everyone rates as 10 points or an item that everyone rates as 3 points, are preferable.

% \paragraph{項目の選出}さてこうして尺度値が計算できたら，それを順に並べていきます。「夜も眠れない」は10.7点，「事務所に行くことがある」は9.5点，「ニュースをなるべく見る」は7.9点・・・というようにしていくことができますね\footnote{中央値なのになぜ小数点があるのだ，と思う人がいるかもしれません。平均値でもいいですし，評定者が偶数人の場合は中央値も両得点の平均や重みつき平均で小数点が出るからです。}。このとき，項目間の間隔が均等になるように項目を選別します。たとえば「夜も眠れない」と「事務所に行くことがある」の間隔は1.2点ですが，「事務所に行くことがある」と「ニュースをなるべく見る」の間隔は1.6点になっています。これでは等間隔と言えないので，1.2点間隔すなわち8.3点ぐらいの項目を選出します。
\paragraph{Selecting Items} Now that we can calculate the scale values, we line them up in order. "Can't sleep at night" is 10.7 points, "Going to the office" is 9.5 points, "Watching the news as much as possible" is 7.9 points... and so on\footnote{Some people may wonder why there is a decimal point in the median value. It could be an average and also in the case of an even number of raters, the median can be the average or a weighted average of both scores, hence the decimal points.}. At this point, select items so that the intervals between each are consistent. For instance, the interval between "Can't sleep at night" and "Going to the office" is 1.2 points, but the interval between "Going to the office" and "Watching the news as much as possible" is 1.6 points. This is not an equal interval, so we select items around 8.3 points, which corresponds to a 1.2 point interval.


% このことからわかるように，サーストン法で尺度を作る場合は，事前に多くの項目を準備しておかないと「ちょうどいいところの表明文がない」となってしまう恐れがあります。ですから最終的にできる尺度に含まれる項目の，5倍から10倍ぐらいを事前に準備し，うまく等間隔に項目が選出できるようにしなければなりません。
As can be seen from this, when creating a scale using the Thurstone method, there is a risk of ending up with "no statement in the right place" unless many items are prepared in advance. Therefore, about 5 to 10 times the number of items included in the final scale must be prepared in advance, and you must ensure that items can be selected at equal intervals.


% なぜ等間隔に選ぶのかというと，もうお分かりですね，これで得られる尺度値を\keytermL{間隔尺度水準}{かんかくしゃくどすいじゅん}として扱いたいからです。間隔が等しくなければ順序尺度にしかなりませんが，間隔が等しいことがわかっていると，平均や分散などの計算をし，相対的な比較をできるからです。またこの尺度を使うときは，すでに評定者集団によって尺度値がわかっていますから，回答者にずらりと並べられた尺度を見てもっとも自分の意見に近い項目を選出してもらえば，その項目の尺度値がその人の態度得点だということができます。
The reason we choose at equal intervals is, as you might already know, that we want to treat the obtained scale values at the level of "interval scales". If the intervals are not equal, they will only become ordinal scales, but if we know that the intervals are equal, we can calculate averages and variances, and make relative comparisons. Also, when using this scale, we already know the scale values by the raters, so if the respondents can choose the item that is closest to their opinion from the line-up of scales, we can say that the scale value of that item is the score of that person's attitude.

% 評定者集団をなるべく偏りなく多く集めることで，事前に尺度の値を確定させておき，あとは本来研究対象にしたかった人にその尺度を当てれば尺度値(尺度得点)が求められる方法ですから，準備が大変だけど使うときは確実で絶対的なスコアを与えることができるというのがこの方法の利点です。欠点はその準備コストの高さと，1次元的な態度しか用いられないことでしょうか。また尺度構成の観点から重要なのは，選出プロセスによって項目間の尺度値が均等であることが保証されている点です。均等に選んだ後で，1，2，3，4，5と数字を付け直しても構いません。大事なのは，こうしたプロセスのおかげで間隔尺度水準が維持され，以後の分析に耐えうるスコアになっているという点です。
By gathering as many rater groups as possible without bias, the value of the scale is predetermined, and then you can get the scale value (scale score) by applying that scale to the person you originally wanted to study. The advantage of this method is that it gives a reliable and absolute score, though preparation is difficult. The drawbacks are the high cost of preparation and the fact that it can only use a one-dimensional attitude. Another important point from the perspective of scale construction is that the scale values between items are ensured to be equal by the selection process. After choosing evenly, you can reassign numbers 1, 2, 3, 4, 5. What's important is that thanks to this process, the interval scale level is maintained and the score can withstand subsequent analysis.

% \section{リッカートのシグマ法}
\section{Likert's Sigma Method}
\label{LikertScale}
% 次に紹介するのはリッカートの\keyterm{シグマ法}{しぐまほう}{$\sigma$ method}です。これはいわゆる5件法，7件法と呼ばれる採点方法で，「私は自民党の政治のやり方が好きだ」といった項目に対して，「まったく当てはまる」「かなり当てはまる」「やや当てはまる」「どちらとも言えない」「やや当てはまらない」「あまり当てはまらない」「まったく当てはまらない」といった順序づけられたカテゴリーにたいしてもっとも自分の考え・態度と近いところに丸をする，という方法で反応が得られます。この時の反応カテゴリーが，今回は7つありますから7件法(7-points scale)で回答を求めた，などと言います。5段階なら5件法，4段階なら4件法です。普通は「どちらとも言えない」というところを用意するために奇数(3,5,7,9)件法を使いますが，日本人は「どちらとも言えない」を選びやすいという中庸傾向があるとも言われていますので，意見をはっきりさせるために4,6件法も使われたりします。
The next introduction is Rikert's \keyterm{Sigma method}. This is a scoring method commonly referred to as a 5-point scale or 7-point scale. For the statement, "I like the way the Liberal Democratic Party operates politically," respondents select and mark a category that is closest to their own thoughts or attitudes from ordered categories such as "completely agree," "largely agree," "somewhat agree," "neither agree nor disagree," "somewhat disagree," "largely disagree," and "completely disagree." The number of response categories here is seven, so it's referred to as 7-point scale. If it's five stages, then it's a 5-point scale, and for four stages, it's a 4-point scale. Scales with an odd number of choices (3, 5, 7, 9) are commonly used to provide an option for "neither agree nor disagree," but it is said that there is a tendency for Japanese to easily choose "neither agree nor disagree," so 4-point or 6-point scales are sometimes used to make opinions clearer.

% 項目はこれも複数あって，たくさん集められた項目を分析するために，もっとも当てはまるを7，かなり当てはまるを6，以下同様にしてまったく当てはまらないを1，とコード化し分析するのが一般的です。ただし注意して欲しいのは，もっとも当てはまる＝7としたのは\keytermL{名義尺度水準}{めいぎしゃくどすいじゅん}の数字の割り当て方と一緒で，このカテゴリーが7という尺度値を持っているわけではない点です。そもそもこの評定カテゴリーは，統計学的にはせいぜい順序尺度水準の性質しか持っていませんから(もっとも＞かなり＞やや)，カテゴリーに割り当てた数字からそのまま平均や分散の計算をするのはおかしいはずなのです。もしこれらのカテゴリーの間隔が等しかったとしても，3，4段階しかないようであればやはり間隔尺度水準の計算ができるほどの精度は持っていません。数量的に分析するには(等間隔が担保された上で)9から11段階は必要と言われています。
There are also multiple items, and to analyze the collected items, it is common to code and analyze by matching the most applicable as 7, quite applicable as 6, and so on, with completely inapplicable as 1. However, what I want you to pay attention to is that the fact that the most applicable = 7 is the same as the way numbers are assigned at the nominal scale level, this category does not necessarily have a scale value of 7. After all, this rating category has at most the property of an ordinal scale level in statistics (most > quite > somewhat), so it should be strange to calculate the average and variance directly from the numbers assigned to the categories. Even if the intervals between these categories were equal, if there were only 3 or 4 levels, they would not have the precision to perform interval scale level calculations. To analyze quantitatively (assuming equal intervals), it is said that 9 to 11 levels are necessary.

% それではリッカート法ではどのようにして尺度値を決めるのでしょうか。リッカート法も測定しようとしているのは態度であって，表に出てくる反応カテゴリーの背後には連続的な心理的態度というのがある，と仮定しています。またこの(社会)心理学的態度は，向きと大きさがあって正規分布を仮定できます。リッカート法も正規分布に従う潜在的な連続変数があると仮定するのです。
So how do you determine the scale values in the Likert method? The Likert method also assumes that what it is trying to measure is attitudes, and there are continuous psychological attitudes behind the response categories that come out in the table. It also assumes that this (social) psychological attitude has direction and size and can assume a normal distribution. The Likert method also assumes that there is a latent continuous variable that follows a normal distribution.

% さて，ある項目について，多くの人からデータを集めて「もっとも当てはまる」「かなり当てはまる」といったカテゴリーごとの集計ができたとしましょう。多くの人の態度も集積すれば正規分布に従いますから，きっとこのヒストグラムも正規分布を反映したものになっているはずです。しかし我々が知りたいのはその背後にある連続体上の数字なわけです。
Now, suppose we have collected data on a certain item from many people and can aggregate it into categories such as "most applicable" and "quite applicable". If we accumulate the attitudes of many people, it will follow a normal distribution, so this histogram should reflect a normal distribution. However, what we want to know is the continuous numbers behind it.

\begin{figure}[ht]
	\centering
	\includegraphics[keepaspectratio,width=0.8\linewidth]{../images/02_Scaling/Rplot02_01.png}
	\caption{カテゴリ反応と背後の連続値}
	\label{fig::t2_02_01}
\end{figure}


% ここで図\ref{fig::t2_02_01}を見てください。上段にあるのがある項目のヒストグラムの例です。しかし知りたいのは，下段にあるような正規分布の形をした連続体の変数のはずです。上段のヒストグラムは下段の状態を反映しているはずですから，上段のカテゴリの相対頻度を元に，下段の正規分布を分割します。具体的な数字との対応は表\ref{tbl::02_01}を見てください。出現度数を相対頻度にし，正規分布の面積を順に分割していくことになります。カテゴリの下の方から順に分割するということで，表\ref{tbl::02_01}の三段目には累積(相対)頻度を書いてあります。そしてこれを使って，標準正規分布の下から面積を考えます。統計環境Rでは，\texttt{qnorm}関数をつかうと累積確率の確率点が求められるのでした(表\ref{tbl::02_01}の4段目)。またその時の確率密度も求めてあります(表の5段目。Rでは\texttt{dnorm}関数を使って求めます)。
Please refer to Figure \ref{fig::t2_02_01}. The histogram of a certain item is shown on the top row. However, what we want to know should be continuous variables shaped like a normal distribution, as shown in the bottom row. Since the histogram in the top row should reflect the state in the bottom row, we divide the normal distribution in the bottom row based on the relative frequency of the categories in the top row. For the correspondence with specific numbers, please refer to Table \ref{tbl::02_01}. Frequencies are converted to relative frequencies and the area of the normal distribution is divided in order. This means dividing from the lower categories, and the cumulative (relative) frequency is written in the third row of Table \ref{tbl::02_01}. And using this, we consider the area from the bottom of the standard normal distribution. In the statistical environment R, the qnorm function was used to find the cumulative probability point (fourth column of Table \ref{tbl::02_01}). Also, the probability density at that time is calculated (fifth column of the table. In R, this is calculated using the dnorm function).

% latex table generated in R 4.0.3 by xtable 1.8-4 package
% Wed Mar 10 11:36:18 2021
% latex table generated in R 4.0.3 by xtable 1.8-4 package
% Wed Mar 10 12:01:46 2021
\begin{table}[ht]
   \centering
   \caption{カテゴリと数値の対応}
   \begin{tabular}{rp{1cm}p{1cm}p{1cm}p{1cm}p{1cm}p{1cm}p{1cm}}
      \hline
      反応カテゴリ      & \pbox<t>{\small{まったく当てはまらない}} & \pbox<t>{\small{あまり当てはまらない}} & \pbox<t>{\small{やや当てはまらない}} & \pbox<t>{\small{どちらとも言えない}} & \pbox<t>{\small{やや当てはまる}} & \pbox<t>{\small{かなり当てはまる}} & \pbox<t>{\small{まったく当てはまる}} \\
      \hline
      出現度数        & 16                            & 46                           & 111                         & 142                         & 126                       & 42                         & 17                          \\
      相対度数        & 0.03                          & 0.09                         & 0.22                        & 0.28                        & 0.25                      & 0.08                       & 0.03                        \\
      累積相対度数      & 0.03                          & 0.12                         & 0.35                        & 0.63                        & 0.88                      & 0.97                       & 1.00                        \\
      累積相対度数の確率点  & -1.85                         & -1.16                        & -0.40                       & 0.33                        & 1.19                      & 1.83                       & $\infty$                    \\
      累積相対度数の確率密度 & 0.07                          & 0.20                         & 0.37                        & 0.38                        & 0.20                      & 0.08                       & 0.00                        \\
      付与される尺度値    & -2.33                         & -1.44                        & -0.77                       & -0.04                       & 0.72                      & 1.50                       & 2.67                        \\
      \hline
   \end{tabular}
   \label{tbl::02_01}
\end{table}


% さて，ではここからどのようにして尺度値を求めればいいでしょうか。一般にC件法で，下から$1,2,3,\cdots c\cdots, C$とカテゴリ順に数字を割り振ったとして，第$c$カテゴリの尺度値$Z_c$を考えるとします。このカテゴリ$c$は標準正規分布において上限$z_c$，下限$z_{c-1}$の確率点で挟まれる領域としていますから，この幅([$z_{c-1} , z_c$])の平均を取ることを考えます。この点は，次の式で求められます(証明は付録\ref{ProofOfLikertScore},Pp.\pageref{ProofOfLikertScore}参照)。
Now, how should we derive the scale value from here? In general, with the CC method, we assign numbers to categories in ascending order from $1,2,3,\cdots c\cdots, C$, considering the scale value $Z_c$ for the $c$th category. Category $c$ is considered as a region trapped by probability points with a lower limit of $z_{c-1}$ and an upper limit of $z_c$ in the standard normal distribution. Now, consider taking the average of this width (between [$z_{c-1} , z_c$]). This point can be calculated by the following formula (refer to Appendix \ref{ProofOfLikertScore}, Pp. \pageref{ProofOfLikertScore} for the proof).

\[ Z_c = \frac{(y_{z_{c-1}} - y_{z_c})}{p_c} \]

ここで$y_c$は$z_c,z_{c-1}$における確率密度，$p_c$はカテゴリ$c$の相対頻度です。具体例でいきましょう。表\ref{tbl::02_01}の数字を使うと，「やや当てはまらない」($c=3$)の尺度値は，$\frac{0.20-(0.37)}{0.22}=-0.7727273$となります(分子は図\ref{fig::t2_02_01}の点線部の差分，分母は該当箇所の面積になります)\footnote{表\ref{tbl::02_01}は小数点下2桁までに丸めているので，正確な値ではありません。}。このようにして計算された尺度値が，表\ref{tbl::02_01}の一番下の段にある数字です。

このように，累積度数をつかって尺度値を決めるリッカートの方法を\keytermL{シグマ法}{しぐまほう}と言います。こうして作られた尺度値は連続体上の数字ですから間隔尺度水準になり，平均や分散をはじめとした数値計算に耐えうる値になっています。機械的に「まったく当てはまらない」から「まったく当てはまる」まで，$1.0$刻みで数字を割り振っているのではないのです！


・・・と言いたいところですが，今回の尺度値を眺めてみるとそこそこ等間隔(間隔は$0.8$ぐらいでしょうか)に並んでいますね。試しに各尺度値を0.8で割ってみると，$-2.91,-1.81,-0.97, -0.04,0.90, 1.88, 3.33$となります。四捨五入して小数点をなくしてみると，$-3,-2,-1,0,1,2,3$となりますね。そう，つまり\textbf{非常にラフな近似値でよければ，機械的に1,2,3...と数字を振っても同じことになります}。ですから，実際の研究ではとくに深く考えずに1,2,3...と割り振った数字をそのまま使われたりするのです。
大山鳴動して鼠一匹といいますか，苦労した割に得るものがないじゃないか，とお叱りを受けそうですが，少なくとも「なぜリッカート法は順序尺度ではなく間隔尺度のように扱って良いのか」という問いには答えられると思います。また，ここにくるまでに，態度の1次元性や正規分布の仮定などが含まれていたことを改めて思い出してください。今回は数値例ですので，綺麗に七段階に分かれるようなものを用意しましたが，実際の調査では正規分布しないものや，平均が低すぎるとか高すぎるものが結構みられます。それらに対して機械的につけた数字で分析するのは決して適切な方法ではなく，シグマ法を始めその他の手法で適切な尺度値を付与すべきなのですが，人間は易きに流れるもので\textbf{ほとんど考慮されていないのが現状です}\footnote{この状況は決して良いものではなく，悪しき研究上の風習だと思われます。幸い，第\ref{x04_IRT}講で説明する\keyterm{項目反応理論}{こうもくはんのうりろん}{Item Response Theory}の一種，\keyterm{段階反応モデル}{だんかいはんのうもでる}{Greaded Response Model}を用いると，この問題点をカヴァーしつつ有用な情報が得られますので，みなさんは一足飛びにその手法を身につけた方が良いと思います。}。

\section{尺度を評価する}
このようにして作られた心理尺度は，それがきちんと測定したいものを測定できているか，評価する必要があります。ここでは\keyterm{信頼性}{しんらいせい}{reliability}と\keyterm{妥当性}{だとうせい}{validity}の2つの側面から説明します。
\subsection{信頼性}
信頼性は測定の安定性と言い換えてもいいかもしれません。すなわち，同じものを2回測っても同じ数字がつくことですね。測定するたびに数字が変わるようでは，その測定器(ここでは尺度ですが)は信用ならない，というわけです。

テスト理論の文脈では，テストのスコア$X$を本当に測りたいもののスコア$t$と誤差$e$とに分割して考えます。古典的テスト理論のモデル式は次の通りです。

\[ X = t + e \]

% ここで，各項目についても同じことが言えると考えると，$X_i = t_i + e_i$ということになります。複数項目で測定するのは，これの平均値を考えると$\bar{e}=0$(誤差の平均値はゼロになる)という仮定から，$\bar{X}=\bar{t}$となって真のスコアを得ることができる，ということがわかります\footnote{この点については，心理学データ解析基礎の授業でも触れていますので，あっさりとした説明になっています。よくわからない人は「基礎」の方の資料を読み直して確認しておいてください。}
Here, the same can be said for each item, which would be $X_i = t_i + e_i$. Measuring multiple items suggests that considering the average value of this results in $\bar{e}=0$ (the average error value is zero), from which we get $\bar{X}=\bar{t}$ and we can obtain the true score, as you can see\footnote{This point is also touched on in the basic psychology data analysis class, so the explanation is very straightforward. If you don't understand, please go back and review the "basics" materials for confirmation.}
% また，テストの分散$Var(X)$を考えると，テスト全体の分散は$Var(X)=Var(t)+Var(e)$ となり，真のスコアの分散と誤差の分散に分解できることがわかります。ここから，信頼性$Rel$を次のように表現できます。
Also, when considering the variance $Var(X)$ of the test, it is clear that the total variance of the test can be broken down into the variance of the true scores and the variance of the errors, $Var(X)=Var(t)+Var(e)$. From here, reliability $Rel$ can be expressed as follows.

\[Rel = \frac{Var(t)}{Var(X)}= \frac{Var(X)-Var(e)}{Var(X)} = 1-\frac{Var(e)}{Var(X)}\]

言葉で言えば，信頼性は全分散に占める真分散の割合であり，全体から誤差分散の割合を引き算したものであるとも言えます。

信頼性のない尺度があれば，その後のお話は先に進みませんから，まずもって「尺度が信頼できるかどうか」を評価する必要があります。このことを\underline{信頼性は妥当性の上限}，と表現したりします。
信頼性を評価する方法は，測定値の安定の程度を評価できればいいのですから，複数の測定を持ってその相関係数を計算することでひとまず達成できます。
しかし同じ尺度を何度も使うのは，調査回答者に要らぬ構えを持たせてしまいますから，普通は1回の尺度を分割してその特徴を見ることにします。
尺度全体から計算される回答者の値は，項目の尺度値の合計であるのが普通です(テストの点数も正答数に対応していますね)。ですから，ある項目$j$の尺度値は，$j$を除いた残り$M-j$個の尺度値の和と高い相関をするはずです。このように，各項目が尺度全体の値とどの程度相関しているかを見る\keyterm{IT相関}{ITそうかん}{Item-Total correlations}は，尺度の信頼性を見る1つの指標になります。ある項目が，尺度全体と相関していなければ，それはその項目が尺度の中で目的と違うものを測定している可能性があり，それは必然的に尺度の安定を損ねる結果になるからです。

この考え方を発展させ，各項目が他の項目とどの程度相関しあっているか，つまり尺度の中でどの程度整合性がとれたもの＝一貫して同じものを測定しているのかを評価する指標として，\keyterm{$\alpha$係数}{あるふぁけいすう}{alpha coefficient}があります\footnote{\keyterm{クロンバックのアルファ}{くろんばっくのあるふぁ}{Cronbach's alpha}とも呼ばれます。}。これは，テストに含まれる項目数をM，テスト全体の分散を$V_t$，項目jの分散を$V_j$と表すと，次の式で表されます。

\[  \alpha =   \frac{M}{M-1} \times \left(1-\frac{\sum V_j}{V_t} \right) \]

% この指標は，各項目が他の項目とどれぐらい相関するかを総合的に表した指標で，とくに\keytermL{内的整合性信頼性}{ないてきせいごうせいしんらいせい}と呼ばれます。
This indicator is a comprehensive representation of how much each item correlates with other items, particularly known as Internal Consistency Reliability.
% このようにして，尺度の安定の程度である信頼性は数値化できますが，次にお話しする妥当性については，数値化できないものです。
In this way, reliability, which is the degree of stability of the scale, can be quantified, but validity, which we will talk about next, cannot be quantified.
% \subsection{妥当性}
\subsection{Validity}
% \keyterm{妥当性}{だとうせい}{validity}は，信頼性をその上限とした上で，それが何を測っているのかを改めて考える指標です。
"Validity", is a measure that considers what it is measuring, with reliability as its upper limit.

% たとえば身長を測ろうとして，体重計を使うとします。成長に応じて，身長が伸びますが，それは体重とも関係がありますので，身長の伸びに応じて体重も増えていくでしょう。体重計は安定した計測器で，信頼性は十分あると思いますが，体重計で身長が測れていると言えるでしょうか。相関する変数ですので，部分的にYesといえそうですが，やはり身長は身長計で測ったほうが良いでしょう。身長計の方が，身長という特徴を的確に捉え，より本質に近い測定をしているからです。このように，作ったものがしっかりとその本質を捉えているかどうか，これが妥当性の基本的なポイントです。
For instance, let's say you are trying to measure your height with a weighing scale. As you grow, your height increases, but so does your weight, as it is correlated with your height. The weighing scale is a consistent measurement tool, and it's fair to say it's pretty reliable. But can we say that it measures height? Since these variables are correlated, we could partially say yes; however, it's always better to measure height with a height measuring scale. The reason being is that a height measuring scale more accurately captures the characteristic of height, providing a measurement closer to the essence of the matter. This is a fundamental point of validity - whether or not what we've created accurately captures its essence.

% 妥当性はですから，そもそも概念としてその測定しようとしているものが適切かどうか(\keyterm{構成概念妥当性}{こうせいがいねんだとうせい}{construct validity})とか，その文言でちゃんと質問できているか(\keyterm{内容的妥当性}{ないようてきだとうせい}{content validity})，理屈通りその測定値が結果と変動しているか(\keyterm{基準関連妥当性}{きじゅんかんれんだとうせい}{criterion validity})と言った面から検証されます。最後の基準関連妥当性については，基準値と尺度値をつかって数量的に検証できますが，構成概念妥当性や内容的妥当性は中身の問題であったり，言葉と概念の対応であったりするので，数理モデル的アプローチができるものではありません。数値化できないか大きな問題ではない，というのはもちろん逆で，数値化できないところであるからこそ，専門的な観点，幅広い視野，批判的思考でもって検証していかなければなりません。
Validity is examined from various aspects such as whether the concept it is trying to measure is appropriate in the first place (construct validity), whether it can properly ask questions with its wording (content validity), and whether the measurements vary as logically expected (criterion validity). For the last criterion validity, it is possible to verify quantitatively using the criterion and scale values, but as for the construct validity and content validity, it a problem of content or a correspondence between words and concepts, so it is not something that can be approached with mathematical models. The fact that it cannot be quantified is not a big problem, on the contrary, precisely because it cannot be quantified, it must be verified with professional viewpoints, a wide field of view and critical thinking.

% 量的に評価する方法としては，今後説明していく\keyterm{因子分析}{いんしぶんせき}{Factor Analysis}によって\keyterm{因子的妥当性}{いんしてきだとうせい}{factorial validity}を見る方法ですとか，理論通りの因子に分離できているかを見る\keyterm{弁別的妥当性}{べんべつてきだとうせい}{distinctive validity}，あるいは\keyterm{収束的妥当性}{しゅうそくてきだとうせい}{convergent validity}などがあります。これらをまとめて，\keyterm{検証的因子分析}{けんしょうてきいんしぶんせき}{confirmatory factor analysis}によって理論通りの分類ができているかをモデル適合度の観点から評価する方法もあります。これらは次回以降お話ししていく，テスト理論の発展系のなかで考えていくものですので，どうぞお楽しみに。
As a quantitative evaluation method, there is a method to see the factorial validity with factor analysis, which we will explain later, the distinctive validity to see if it can be separated into factors as theorized, or the convergent validity, etc. These are collectively evaluated from the perspective of model fitness by confirmatory factor analysis to see if the classification is as theorized. These are what we will consider in the development of test theory, which we will discuss from the next time on, so please look forward to it.

% \section{課題}
\section{Tasks}
% \paragraph{リッカート法}シグマ法でなく機械的に数字を割り振るとどのような問題が生じるか，自分で数値例を作って検証してみてください。
\paragraph{Likert Scale} Please create a numerical example yourself and check what problems arise when assigning numbers mechanically, not using the sigma method.
% \paragraph{信頼性の記述と報告}心理学の尺度作成に関する論文\footnote{日本心理学会が出している「心理学研究」という学会誌では，【資料】というカテゴリーで毎回のように新しい尺度が報告されています。}を読み，信頼性についてどのように記述されているかを確認してみましょう。
\paragraph{Description and Reporting of Reliability} Let's check how reliability is described by reading papers on scale creation in psychology\footnote{In the academic journal "Psychological Research" issued by the Japanese Psychological Association, new scales are reported almost every time in the category of "Materials"}.
% \paragraph{さまざまな妥当性}妥当性にはさまざまなものがあります。\textcite{kosugi2016b}などを参考に，妥当性について自分なりに調べてみてください。
\paragraph{Various Validities} There are various kinds of validity. Please research about validity on your own, referring to sources like \textcite{kosugi2016b}.
% \paragraph{リッカートのシグマ法} 表\ref{example2_1}のように，「かなり当てはまる」や「まったく当てはまる」など尺度の右の方に丸をつける人が多かった項目があったとします。この時の尺度値をリッカートのシグマ法に則って算出してみてください。
\paragraph{Likert's Sigma Method} Suppose there was an item where many people marked the right side of the scale, such as "quite applicable" or "completely applicable", as shown in Table \ref{example2_1}. Please calculate the scale value according to Likert's Sigma method at this time.

\begin{table}[ht]
	\centering
	\caption{天井効果の出た尺度}
	\begin{tabular}{rp{1cm}p{1cm}p{1cm}p{1cm}p{1cm}p{1cm}p{1cm}}
		\hline
		反応カテゴリ      & \pbox<t>{\small{まったく当てはまらない}} & \pbox<t>{\small{あまり当てはまらない}} & \pbox<t>{\small{やや当てはまらない}} & \pbox<t>{\small{どちらとも言えない}} & \pbox<t>{\small{やや当てはまる}} & \pbox<t>{\small{かなり当てはまる}} & \pbox<t>{\small{まったく当てはまる}} \\
		\hline
		出現度数        & 1.00                          & 3.00                         & 5.00                        & 18.00                       & 24.00                     & 53.00                      & 56.00                       \\
		相対度数        & 0.01                          & 0.02                         & 0.03                        & 0.11                        & 0.15                      & 0.33                       & 0.35                        \\
		累積相対度数      & 0.01                          & 0.03                         & 0.06                        & 0.17                        & 0.32                      & 0.65                       & 1.00                        \\
		累積相対度数の確率点  & -2.50                         & -1.96                        & -1.59                       & -0.96                       & -0.47                     & 0.39                       & $\infty$                    \\
		累積相対度数の確率密度 & 0.02                          & 0.06                         & 0.11                        & 0.25                        & 0.36                      & 0.37                       & 0.00                        \\
		\hline
	\end{tabular}
	\label{example2_1}
\end{table}


\end{document}
